\section{Introduction}
\label{sec:intro}

Data scarcity remains a fundamental bottleneck in robotics, motivating the use of inexpensive and abundant surrogate data such as simulation and cross-embodiment data~\cite{bjorck2025gr00t, intelligence2025pi}. 
Although these data sources contain rich task-relevant information, they introduce substantial domain gaps that make effective knowledge transfer difficult in practice. 
Recently, a simple \textit{co-training} paradigm — jointly training on in-domain real data and surrogate data with a data mixing ratio $w$ — has demonstrated strong empirical performance across sim-and-real and human-to-robot settings~\cite{cheng2025generalizable, yuan2025motiontrans, kareer2025egomimic, kareer2025emergence}. 
Despite some work~\cite{wei2025empirical} providing valuable empirical analysis, co-training remains poorly understood: its internal mechanisms are largely treated as a black box, and the factors that govern its effectiveness are unclear.

In this work, we focus on investigating the co-training paradigm as described above, in particular, on \textit{sim-and-real} data, in diffusion-based models, which are representative in modern generative robot policies.~\cite{pan2025much}



We begin with a theoretical analysis that examines the learning objective induced by jointly mixing data from multiple domains.
This analysis reveals two intrinsic effects that independently influence co-training performance:
% \textbf{structured representation alignment} and \textbf{softmax lensing effect}.
(1) \textbf{\emph{Structured representation alignment}} is characterized by a two-fold property. On one hand, representations become aligned across domains in a domain-invariant subspace, enabling the transfer of task-relevant knowledge. On the other hand, representations retain discernibility with respect to domain-specific factors, allowing actions to adapt to the real world rather than being directly copied from surrogate domains.
This balance is instrumental for effective co-training, as it determines whether adaptive action transfer is possible.
(2) \textbf{\emph{Importance reweighting effect}} refers to the domain-dependent logit modulation within action weightings.
This effect operates locally in the space conditioned on observations and controls how much each training sample from each domain contributes to an action decision during training.
It is determined by the data mixing ratio $w$, the dataset size $|\mathcal{D}|$ and domain gaps.
Through controlled toy co-training experiments, we verify the presence of both effects and find that structured representation alignment is the primary factor underlying strong model performance, while the importance reweighting effect plays a modulatory role.
These findings motivate the following questions:

\textit{
Do similar effects arise in realistic robot manipulation tasks?
How can these insights guide the design of more effective co-training algorithms?
}




In the second part of the paper, we address these questions through comprehensive sim-and-sim and sim-and-real robotic manipulation experiments.
In end-to-end co-training systems, the data mixing ratio $w$ is typically the only explicit control variable, yet it simultaneously influences both internal effects.
Empirically, we find that structured representation alignment, in both local and global space, can emerge implicitly within an appropriate range of mixing ratios (which we refer to as ``balanced mixing ratios"), and that its strength exhibits a moderate-to-strong correlation with task success.
At the same time, preserving domain discernibility is necessary for effective action adaptation to the real world; when this property is lost, performance even shows a negative correlation with representation alignment.
These observations provide a unified perspective for understanding existing co-training techniques. We benchmark three recent representative co-training techniques on our tasks --- optimal transport-based feature regularization~\cite{cheng2025generalizable, punamiya2025egobridge}, adversarial discriminative domain adaptation~\cite{cai2025n, yuan2025motiontrans}, and classifier-free guidance~\cite{wei2025empirical}. 
We observe that each method primarily emphasizes only one aspect of structured representation alignment, which often leads to unstable or marginal improvements.
Motivated by this analysis, we propose a simple combination of co-training techniques that jointly promotes alignment while preserving domain discernibility, and that also offers a more controllable interface for knowledge transfer during inference. This approach yields consistent and substantial improvements over prior methods.
 In summary, our contributions are as follows.
\begin{itemize}[leftmargin=*]
    \item We systematically identify, for the first time, the working mechanisms of co-training through theoretical analysis and experimental support.
    \item We find that structured representation alignment can be learned implicitly, and validate the effects and requirements of both alignment and discernibility through comprehensive robotic manipulation experiments.
    \item We benchmark representative co-training techniques through the lens of our analysis, which inspires us to introduce a simple approach that stably improves performance, and opens the door to new algorithm designs.
\end{itemize}







